% 第六章 討論
\chapter{討論}\label{ch:discussion}

\section{適用條件}\label{sec:when}

尾數歸零的效益取決於以下三個條件，缺少其中任何一個，收益就會縮小或消失：

\begin{enumerate}
  \item \textbf{系統受功耗限制。}省下的功耗只有在功耗上限壓住時脈時才會轉成速度。HPL 在 500–800~W 快 4.0–5.0\%（低於規格表下限的 400~W 為 11.0\%）；調參後的 HPL 在預設的 1000~W 仍有 75\% 的時間撞上限，所以每次求解少耗電 7.5\%，速度也快約 4\%（兩對差異較大）。幾乎不受限的程式（例如 1000~W 的 DF-MP2，只有部分時間撞上限）則主要表現為省電，速度只多 0.3–0.8\%。
  \item \textbf{被截斷的資料直接且大量地進入乘法器或資料通道。}矩陣乘法的每個輸入都要參與多次乘法，所以效益最大；N 體計算第一步就做減法，差值的尾數又是滿的；CG 與 SW4 被截斷的資料只佔記憶體流量的一部分；LAMMPS 的位能表常駐快取，幾乎不經過記憶體。
  \item \textbf{有機制把誤差攔下來。}HPL 的殘差檢查、HPL-MxP 的迭代修正與 SCF 的自洽迭代都會修正或偵測截斷誤差；AI 的權重直接決定輸出品質，截斷便直接轉成精度損失。
\end{enumerate}

表~\ref{tab:summary} 彙整各類工作負載在功耗受限時的收益與主要原因。

\begin{table}[tbp]
\centering
\caption{各類工作負載在功耗受限時的收益}\label{tab:summary}
\small
\begin{tabularx}{\textwidth}{@{}>{\raggedright\arraybackslash}p{3.6cm}LL@{}}
\toprule
工作負載 & 收益 & 主要原因 \\
\midrule
具殘差檢查的稠密求解（HPL、HPL-MxP） & HPL 在 500–800~W 快 4.0–5.0\%；HPL-MxP 在 500~W 快 6.5\%、600–800~W 約 2\%（對角優勢測試矩陣） & 矩陣乘法主導，殘差檢查把關 \\
高功耗的量子化學（DF-MP2） & 快 4.9–8.3\%；1000~W 時省電 8–14\% & 矩陣乘法佔 78\%，張量建好後截斷一次即可 \\
跨度寬且功耗高（高斯過程地統計） & 快 7.4–8.2\%；1000~W 時快 0–4\%（依設定）、省電 7–15\% & Cholesky 以矩陣乘法為主；FP32 不可用、模擬較慢 \\
自洽迭代（Quantum ESPRESSO） & 矽快 2.4\%；金表面快 4.0\%（約一半來自迭代減少）；總能量印出的每一位都不變 & 截斷誤差低於當圈的容許誤差 \\
高斯基底量子化學（PySCF） & 預設格點估計約 2\%（400~W 約 4\%）；細格點的纈胺黴素在 450–550~W 實測 3.3–3.5\% & 矩陣乘法只佔總時間 32\% \\
記憶體受限、係數佔流量多（MFEM、SW4） & 快 0.3–2.0\% & 其餘滿尾數的向量稀釋了效益 \\
係數常駐快取（LAMMPS） & 落在雜訊內 & 被截斷的資料不經過記憶體 \\
已使用 FP8 的 AI 訓練與推論 & 訓練變慢；推論以精度 1:1 換取速度 & FP8 只剩 3 個尾數位元 \\
\bottomrule
\end{tabularx}
\end{table}

\section{與改變格式做法的關係}\label{sec:vs-formats}

尾數歸零與 INT8 模擬 FP64、FP32／3xTF32、FP8 並非互相取代，而是互補。INT8 模擬在所需精度較低（約 28 位以下），或 $n$ 與 $k$ 都達數千以上時較快；截斷則在瘦長矩陣、需要 32 位以上的精度，以及 GEMMul8 這類高功耗模擬遇到嚴格上限時較快。第 \ref{sec:crossover} 節的三參數公式（式~\eqref{eq:crossover}）在測過的範圍內能描述勝負，但只以單一 $m$ 擬合、形狀與節點差異未分開、樣本外案例都離分界很遠，目前只是經驗參考，還不足以直接作為函式庫選擇方法的依據。FP32 與 3xTF32 的誤差隨內積長度成長，把所有 GEMM 換成它們，在長內積的平面波 DFT 上無法收斂；截斷保留 FP64 累加，誤差與內積長度無關。FP8 已經把 AI 的尾數減到 3 位，在其上再截斷只是以精度換取速度。

整個矩陣改用同一種格式的做法還有一個前提：資料的指數跨度必須夠窄。逐張量縮放的格式（FP32、FP8）取決於整體跨度，逐列切片的 INT8 模擬取決於同一列內的跨度（第 \ref{sec:span} 節）。DF-MP2 的張量跨度窄，所以內建的 FP32 反而比截斷快；高斯基底 DFT 依運算量加權有 24–44\% 的元素（最寬的格點原子軌域矩陣達 87\%）比所在列的最大值小 $2^{-32}$ 以上，固定位數的模擬便在沒有警告的情況下算錯。不過列內跨度寬只是必要條件（高斯過程的跨度更寬卻沒有算錯），以 tile 為單位選擇格式的混合精度也能處理這類資料。這個跨度可以用攔截工具事先量測，作為風險指標；cuBLAS 的自動模式也會偵測它而改用原生 FP64，但如此便不再比截斷快（第 \ref{sec:span} 節）。

因此，尾數歸零的定位是：在 FP64 科學計算中，一個每個元素只多一次位元運算、精度可以 1 位為單位調整、誤差不受矩陣形狀與指數跨度影響的預設選項。在跨度窄、矩陣大而方、精度需求較低時，改用模擬或 FP32 可能更快；對於能依空間排序把小數值集中到整個 tile 的應用（例如地統計），tile 層級的混合精度能快得更多，截斷則可以疊加在它剩下的 FP64 tile 上。

\section{限制與效度威脅}\label{sec:limits}

\begin{enumerate}
  \item \textbf{單一平台，而且是對截斷有利的平台。}所有實驗都在同一個 GH200 叢集上進行。GH200 的 FP64 tensor core 相對強，是 INT8 模擬優勢最小的平台之一：Uchino 等人的 DGEMM 模擬在 RTX 5080（FP64 很弱）上快 18.5 倍，在 A100 與 GH200 上都只快 1.4 倍~\cite{Uchino2025}。在 FP64 相對弱的平台上，截斷與模擬的勝負可能反過來；其他 GPU 世代的功耗管理也不同，效益大小與分界公式的係數都可能改變。GEMMul8 的比較使用 3.5.1 版（2026 年 9 月），之後的版本可能改變結果。
  \item \textbf{保留位數需以多組問題大小驗證。}HPL 的殘差隨矩陣大小跳動，28 位在三組 $N$ 通過（殘差都接近門檻）、在一組失敗；本報告以有百倍餘裕的 36 位作為操作點。
  \item \textbf{應用覆蓋有限。}每個科學軟體只測試一到兩個輸入系統，沒有涵蓋幾何最佳化、長時間分子動力學等情境。
  \item \textbf{攔截工具的複製成本（依應用約 3–9\%）。}因此攔截版本一律與「複製但不截斷」的基準比較，只有程式內版本與未修改的原生程式比較。
  \item \textbf{節點差異。}同一功耗上限下，節點間的 SM 時脈可相差 13–17\%，因此所有比較都在同一組節點上成對進行。
  \item \textbf{散熱條件。}風扇轉速低時 GPU 會先撞到溫度上限；風扇全開時，原生 HPL-MxP 在 1000~W 會觸發機箱的硬體功耗煞車。這些條件都會改變結論，需要一併記錄。
  \item \textbf{共用叢集。}其他使用者的工具至少兩次重設叢集的功耗上限；我們在每次執行前檢查實際生效的上限，受影響的量測一律重做。
  \item \textbf{功耗感測器。}功耗讀自驅動程式，沒有外接電表；在 A100／H100 上，驅動程式只取樣部分時間~\cite{Yang2024}。我們的比較都是同一量測路徑下的成對相對值，但絕對能量仍需以外接電表驗證。模組功耗也包含 Grace CPU 與記憶體（例如 ExaGeoStat 在 CPU 產生共變異數時就與 GPU 共用同一個上限），所以本報告的耗電數字都應視為同一量測路徑下的相對值。
  \item \textbf{訓練的收斂點有雜訊。}LoRA 的達標點即使設定完全相同也會相差一個評估區間，而我們只用了兩個亂數種子，所以只把它當作「FP8 沒有再截斷的空間」的佐證。
  \item \textbf{HPL-MxP 的測試矩陣。}我們使用嚴格對角優勢的矩陣與近似的對角分解，規則只允許這種矩陣用於測試~\cite{HPLMxPRules}；「保留 3 位仍穩定」可能只是這個矩陣的性質。
  \item \textbf{樣本數與節點。}信賴區間只用於矩陣乘法微基準與 Quantum ESPRESSO；HPL、HPL-MxP、DF-MP2 與高斯過程每一格只有兩輪，且在同一組節點上量測。同一工作內兩對的差異遠小於效應（HPL 在 400–800~W 相差最多 0.6 個百分點），但兩對不足以給出信賴區間，也不能排除節點組本身的影響。QE 的暖機是以事後訂的規則排除，應改為在交錯順序之前先執行一次不計入的暖機。分界點掃描的四個 $n$ 值分別在不同節點上量測，形狀效應與節點差異無法分開，是實驗設計上的缺陷。
  \item \textbf{功耗上限的意義。}1000~W 是預設也是最高上限，並不代表不受限：HPL 與 FP64 矩陣乘法在 1000~W 仍有部分時間撞上限。撞上限時 1 秒平均功耗常低於上限，限電器依據的時間尺度不明；低於 450~W 的上限超出規格表的可設範圍~\cite{NVIDIAGH200}。
  \item \textbf{機制實驗的範圍。}填值對照只做了保留 0 位與 20 位，500~W 只在一個節點上量測兩次；保留 0 位時「數值為 0」的額外省電原因尚未驗證。
  \item \textbf{收益不全來自功耗。}在 Quantum ESPRESSO 金表面上，自動調整的收益約一半來自 Davidson 迭代減少（第 \ref{sec:sci-perf} 節），原因尚未查明。
\end{enumerate}

\section{規則合規性}\label{sec:rules}

TOP500 的 HPL 要求以完整的 64 位元浮點運算求解，明文排除以低精度求解再迭代修正，以及以整數運算模擬浮點的方法~\cite{TOP500Linpack}；截斷運算元後的 HPL 也不符合這個要求。HPL-MxP 自 2025 年 6 月起規定最低精度為 16 位元浮點，且嚴格對角優勢的矩陣只能用於測試~\cite{HPLMxPRules}；保留 3 位的 BF16 實質上只有 12 位元，我們的測試矩陣也不符合正式規定。因此本報告所有 HPL／HPL-MxP 數字都是研究結果，並非正式成績。MLPerf Inference 的封閉組允許數學上的量化、禁止剪枝，尾數歸零屬於前者~\cite{MLPerfInfRules}；MLPerf Training 的權重能否截斷需依各項規則判斷，不過就本報告的結果而言，均勻截斷本來就不划算。科學軟體沒有這類限制，判準是物理結果是否在所需的精度之內。
